\sscp{Innovations in the historical phonology}{07-04-01}
We initially discuss the subgroups together,
to compare and contrast.
Of the innovations listed below,
numbers 1-5 could have happened through secondary developments
after splitting from a common parent language.
However, it is difficult to see how numbers 6-9 could have developed
through secondary developments, making a fairly strong case for treating \tsc{Sula-Buru}
and \tsc{\ili{Ambon}-Seram} as separate subgroups with distinct phonological
histories.

\begin{enumerate}
	\item	*-ay and *-aw
		\begin{itemize}[topsep=0pt,leftmargin=\dimexpr 15pt-6pt]	
			\item	\tsc{Sula-Buru}
						and \tsc{Seram-Tanimbar-Bomberai} merge *-aw/*-ay/-*a > **a.
			\item	\tsc{\ili{Ambon}-Seram} has *-ay > \it{a} but retains *-aw = **-aw
						in certain lexical items in \ili{Alune} and related languages.
						See discussion and data in \srf{10-02}.
			\item	\ili{Banda} has *-ay/*-aw > \it{a,} but final *a shows alternations
						of \it{a {\tl} o} which are not attested with reflexes of *-ay/*-aw,
						indicating that we cannot posit straightforward merger of *-ay/*-aw/*-a in \ili{Banda}.
						See  discussion
						and data in \srf{12-05}.
		\end{itemize}
	\item	*mb, *mp, *ma-p
		\begin{itemize}[topsep=0pt,leftmargin=\dimexpr 15pt-6pt]
			\item	\tsc{Sula-Buru} merges *-mp-/*-mb-/*ma-p/*ma-b > **mb,
						but retains the nasal-consonant sequence,
						needed to distinguish them from reflexes of *b = **b
						and *p = **p.
			\item	Nearly all \tsc{\ili{Ambon}-Seram} languages merge *-mp-/*-mb- > **mp (devoiced).
						Data from many languages indicates that *ma-p
						and *ma-b must be kept distinct
						and fully reconstructed for \tsc{\ili{Ambon}-Seram}.
			\item	\tsc{Seram-Tanimbar-Bomberai} keeps *ma-p/*ma-b distinct from *mb/*mp.
						While in some languages they merge,
						in others they have distinct reflexes.
			\item	\ili{Banda} merges *mb/*mp/*ma-p > \it{mb} retaining the cluster
						both initially and medially. 
		\end{itemize}
	\item	*n/*ñ/*ŋ
		\begin{itemize}[topsep=0pt,leftmargin=\dimexpr 15pt-6pt]
			\item	\tsc{Sula-Buru} retains *ŋ = \it{ŋ} root initially
						and medially, but in final position it merges *n/*ŋ > \it{n}{\#}
						in words in which it is retained rather than lost.
			\item	\tsc{\ili{Ambon}-Seram}
						and \ili{Banda} merge *n/*ñ/*ŋ > \it{n} in all positions.
			\item	\tsc{Seram-Tanimbar-Bomberai} reflexes of *ŋ are diverse including:
						\it{ŋ, g, k, ŋg} and \it{n}.
						Notably, \tsc{East Seram} languages (those closest to \tsc{\ili{Ambon}-Seram})
						have medial *ŋ > \it{k}
						and thus do not pattern with \tsc{\ili{Ambon}-Seram} with respect to the nasals.
		\end{itemize}
	\item	*k
		\begin{itemize}[topsep=0pt,leftmargin=\dimexpr 15pt-6pt]
			\item	\tsc{Sula-Buru} regularly retains *k = \it{k} root-initially
						and root-medially, except in final position where it is mostly lost > {\0}.
			\item	\tsc{\ili{Ambon}-Seram} sporadically retain a weakened *k > **k (>
						\it{ʔ}, {\0}), with occasional retentions of initial {\#}\it{k} in lexically
						specific words, particularly verbs that have morphophonemic alternations when inflected.
						(See *nakaw and *kədəŋ in \trf{07-06};
						and examples of such inflections in \trf{07-16}.)
			\item	All branches of \tsc{Seram-Tanimbar-Bomberai} preserve *k =
						\it{k} in at l\ili{east} some words.
						Notably, languages from several branches also have word-final *k = \it{k}.
			\item	\ili{Banda} preserves *k = \it{k} in all word positions in all identified reflexes.
		\end{itemize}
	\item	*q/*h
		\begin{itemize}[topsep=0pt,leftmargin=\dimexpr 15pt-6pt]
			\item	\tsc{Sula-Buru}
						and \ili{Banda} have *q/*h > {\0} in all positions.
			\item	\tsc{\ili{Ambon}-Seram} languages appear to retain conditioned and/or
						lexically specific traces of *q
						and *h in \ili{Kamarian},
						\ili{Alune} and related languages (Naka{\Q}ela,
						\ili{Watui} and \ili{Lisabata}) along with possible traces in others
						\citep[40--52, 133]{Collins1983}.
						See discussion and data in \srf{10-02}, and \srf{13-02-02}.
			\item	Most \tsc{Seram-Tanimbar-Bomberai} languages have *q/*h > {\0},
						but \ili{Watubela} has *q > \it{k} in all positions (\srf{11-03-01}).
		\end{itemize}
	\item	*ə
		\begin{itemize}[topsep=0pt,leftmargin=\dimexpr 15pt-6pt]
			\item	\tsc{Sula-Buru} *ə > \it{e} in both penultimate
						and ultimate syllables (contrasting with \tsc{Taliabu} *ə > \it{o}),
			\item	\tsc{\ili{Ambon}-Seram} languages preserved *ə = **ə with subsequent
						**ə > \it{e, o, i} in complex ways.
						Sometimes this is phonologically conditioned; sometimes it is
						lexically specific at lower level groupings; sometimes it appears
						to be unpatterned, possibly due to contact.
						\citet[65ff]{Collins1983} finds a patterned split in **ə to be insightful
						for subgrouping \tsc{\ili{Piru} Bay} languages.
			\item	\tsc{Seram-Tanimbar-Bomberai} has *ə > \it{a} in final syllables.
						\tsc{Seram-Tanimbar-Bomberai} thus merges *-aw/*-ay/*-ə/*-a >
						**a in the final syllable.
						Most \tsc{Seram-Tanimbar-Bomberai} languages have penultimate
						*ə > \it{e,} though members of \tsc{East Bomberai} show additional complexities.
						See discussion and data in \srf{11-05-02}.
			\item	\ili{Banda} has penultimate *ə > \it{e}
						and a split of final *ə > \it{o,
						e} as conditioned by the penultimate vowel.
						The reflexes of *ə
						and *a in final syllables are distinct in \ili{Banda},
						thus distinguishing \ili{Banda} from \tsc{Seram-Tanimbar-Bomberai}.
						See discussion and data in \srf{12-05}.
			\item[] \hp{\,}
		\end{itemize}
	\item	*z
		\begin{itemize}[topsep=0pt,leftmargin=\dimexpr 15pt-6pt]
			\item	\tsc{Nuclear Sula-Buru} merge *y/*z > **y (> {\0} in \ili{Buru}
						and \ili{Lisela} is a secondary development).
						\ili{Ambelau} has a split of *z > \it{l, h}.
			\item	\tsc{\ili{Ambon}-Seram},
						\tsc{Seram-Tanimbar-Bomberai}, and \ili{Banda} merge *d/*z,
						with **d the likely outcome in \tsc{Seram-Tanimbar-Bomberai},
						but **r in \tsc{\ili{Ambon}-Seram} and \ili{Banda}.
						\tsc{\ili{Ambon}-Seram} has **r/*l > **l in many branches.
						See \srf{07-03-03} above, as well as \srf{10-01} and \srf{10-02}.
		\end{itemize}
	\item	*R/*l
		\begin{itemize}[topsep=0pt,leftmargin=\dimexpr 15pt-6pt]
			\item	\tsc{\ili{Ambon}-Seram} mostly merged *R/*l > **l,
						with complications in Boano
						and \ili{Hoti} meaning that a full merger cannot be posited for those languages.
			\item	In \tsc{Sula-Buru} the distinction is clearly retained with
						*R > \it{h}, *l = \it{l}.
			\item	\tsc{Seram-Tanimbar-Bomberai}
						and \ili{Banda} also retain the distinction (\srf{07-04-01-02}).
		\end{itemize}
	\item	Apicals: *d/*z/*R/*l/*j
		\begin{itemize}[topsep=0pt,leftmargin=\dimexpr 15pt-6pt]
			\item	\tsc{Sula-Buru} strongly preserves the distinctions between
						*d, *z, *R, *l, with *j being complicated (explained below).
			\item	\tsc{\ili{Ambon}-Seram} gives the appearance of merging *d/*z/*R/*l/*j
						> **l, as observed by \citet{Stresemann1927},
						but \citet{Collins1983} argues that they did not fully merge
						in some languages, and *j retains distinct traces,
						so the apparent mergers are probably due to diffusion,
						and such “diffusion has obfuscated differences” (1983:133).
			\item	Different branches of \tsc{Seram-Tanimbar-Bomberai} merge different
						combinations of **d (< *d/*z),
						*R, *l, and *j meaning that four apicals need to be reconstructed.
						The \tsc{Tanimbar-Bomberai} node is defined by *R/*j > **r in
						most (but not all words).
			\item	\ili{Banda} merges *d/*z/*R > \it{r}
						and *j/*l > \it{l} in most words.
		\end{itemize}
\end{enumerate}

Regarding point 9 above,
In a sense \citeauthor{Stresemann1927}
and Collins are both right about the merger of PMP apicals in \tsc{\ili{Ambon}-Seram}.
All 30 \tsc{\ili{Ambon}-Seram} languages for which we have data have
many words which appear to merge *d/*z/*R/*l/*j > **l ---
and a few words that do something else for each of these.
So the appearance of the mergers is quite strong
and quite widespread in \tsc{\ili{Ambon}-Seram} languages,
as noted by \citeauthor{Stresemann1927}.
But the exceptions observed by Collins,
however rare and sporadic, are consistent enough to require some
of the distinctions be preserved for some languages.
Some of the distinctions that need to be preserved in regular
sound correspondences are at the two extremes of \tsc{\ili{Ambon}-Seram},
with Boano in the west (\srf{10-03-01})
and the \tsc{Patakai-Manusela} branch in the \ili{east} (\srf{10-03-03}) being stand-outs.
Sporadic exceptions are also distributed throughout
\tsc{\ili{Ambon}-Seram} (see \tsc{Atamanu} \srf{10-03-02}).
Some of the exceptions are phonologically conditioned,
but many are lexically specific -- some in ways that align with
smaller groupings of languages.
The lexically specific complexity of the reflexes of all these
proto-phonemes can be illustrated with *j.
(See \citet{Collins1983} for a scattered glimpse of other complexities.)

For *j, we see distinct lexically specific patterns: PMP *j >
{\0} in \tsc{Sula-Buru} in *ŋajan ‘name’,
*huaji ‘younger sibling (same sex)’,
but is preserved as *j > **l in \tsc{\ili{Ambon}-Seram} languages in
the same words (except \ili{Kayeli}).
We also see parallel development: e.g.
\ili{Sanana} (\tsc{Sula-Buru}) and South \ili{Wemale} (\tsc{\ili{Ambon}-Seram})
have phonologically conditioned *j > {\0} /Vα{\gap}Vα where both
vowels are the same,
in words in which their related languages normally reflect *j > **l.
This developed independently in these two languages.\footnote{
		Data are limited or patterns are unclear %for subgrouping purposes
		for: *suja ‘pitfall or trail spikes made of sharpened bamboo’,
		*quləj ‘maggot, caterpillar, larva of a metamorphosing insect’,
		*pusəj ‘navel, umbilicus; mid-point’,
		*laləj ‘housefly’, waRəj ‘vine,
		creeper’, *maja ‘dry up, evaporate’,
		*hajək ‘smell, sniff, kiss’, *bujəq ‘foam, bubbles,
		lather, scum, froth’, *qajəŋ ‘charcoal’.}
In addition,
in some languages (e.g.
\tsc{Taliabu}, South \ili{Nuaulu}, \ili{Liana-Seti}) reflexes of *-əj{\#}
are different to reflexes of *ə and *j alone.
Such conditioned reflexes of *-əj{\#} are stable
and different to reflexes of *ə with other consonants.

Similar lexically specific issues and/or phonological conditioning
are found in \tsc{\ili{Ambon}-Seram} languages with PMP *d, *z, *R, *l.
The subtle differences retained in reflexes of *d,
*z, *R, *l, and *j in four branches of \tsc{\ili{Ambon}-Seram} are
discussed in \srf{10-02}. Also see \citet{Collins1983}.